% Options for packages loaded elsewhere
\PassOptionsToPackage{unicode}{hyperref}
\PassOptionsToPackage{hyphens}{url}
\documentclass[
]{article}
\usepackage{xcolor}
\usepackage[margin=1in]{geometry}
\usepackage{amsmath,amssymb}
\setcounter{secnumdepth}{-\maxdimen} % remove section numbering
\usepackage{iftex}
\ifPDFTeX
  \usepackage[T1]{fontenc}
  \usepackage[utf8]{inputenc}
  \usepackage{textcomp} % provide euro and other symbols
\else % if luatex or xetex
  \usepackage{unicode-math} % this also loads fontspec
  \defaultfontfeatures{Scale=MatchLowercase}
  \defaultfontfeatures[\rmfamily]{Ligatures=TeX,Scale=1}
\fi
\usepackage{lmodern}
\ifPDFTeX\else
  % xetex/luatex font selection
\fi
% Use upquote if available, for straight quotes in verbatim environments
\IfFileExists{upquote.sty}{\usepackage{upquote}}{}
\IfFileExists{microtype.sty}{% use microtype if available
  \usepackage[]{microtype}
  \UseMicrotypeSet[protrusion]{basicmath} % disable protrusion for tt fonts
}{}
\makeatletter
\@ifundefined{KOMAClassName}{% if non-KOMA class
  \IfFileExists{parskip.sty}{%
    \usepackage{parskip}
  }{% else
    \setlength{\parindent}{0pt}
    \setlength{\parskip}{6pt plus 2pt minus 1pt}}
}{% if KOMA class
  \KOMAoptions{parskip=half}}
\makeatother
\usepackage{listings}
\newcommand{\passthrough}[1]{#1}
\lstset{defaultdialect=[5.3]Lua}
\lstset{defaultdialect=[x86masm]Assembler}
\usepackage{longtable,booktabs,array}
\usepackage{caption}
\captionsetup[table]{skip=6pt}
\newcounter{none} % for unnumbered tables
\usepackage{calc} % for calculating minipage widths
% Correct order of tables after \paragraph or \subparagraph
\usepackage{etoolbox}
\makeatletter
\patchcmd\longtable{\par}{\if@noskipsec\mbox{}\fi\par}{}{}
\makeatother
% Allow footnotes in longtable head/foot
\IfFileExists{footnotehyper.sty}{\usepackage{footnotehyper}}{\usepackage{footnote}}
\makesavenoteenv{longtable}
\setlength{\emergencystretch}{3em} % prevent overfull lines
\providecommand{\tightlist}{%
  \setlength{\itemsep}{0pt}\setlength{\parskip}{0pt}}
\usepackage{bookmark}
\IfFileExists{xurl.sty}{\usepackage{xurl}}{} % add URL line breaks if available
\urlstyle{same}
% fallback for those not using the hyperref driver hyperxmp:
\makeatletter
\@ifundefined{xmpquote}{\newcommand{\xmpquote}[1]{#1}}{}
\makeatother
\hypersetup{
  hidelinks,
  pdfcreator={LaTeX via pandoc}}

\author{}
\date{}

\begin{document}

\section{OptionPricingApp}\label{optionpricingapp}

\subsection{Multi-Asset Spatiotemporal Stochastic Volatility Options
Engine with Direct Equity Return
Correlations}\label{multi-asset-spatiotemporal-stochastic-volatility-options-engine-with-direct-equity-return-correlations}

This engine simulates a multi-asset basket under a network-coupled
stochastic volatility framework with empirically-calibrated direct
equity return correlations. The model tracks a collection of \(N\)
assets over time \(t \in [0, T]\), where the spot prices are defined by
the state vector \(\vec{S}_t = [S_{1,t}, S_{2,t}, \dots, S_{N,t}]^T\).

\subsubsection{1. Continuous-Time Model Dynamics
(SDEs)}\label{continuous-time-model-dynamics-sdes}

The joint dynamics of the asset prices and their underlying
network-coupled log-variance processes are governed by the following
system of stochastic differential equations:

\[d\vec{S}_t = (\vec{r} - \vec{q}) \odot \vec{S}_t \, dt + \sqrt{\exp(\vec{X}_t)} \odot \vec{S}_t \odot d\vec{W}_t^S\]

\[d\vec{X}_t = \vec{\kappa} \odot \left( \vec{\theta} + \vec{\gamma} \odot (W\vec{X}_t - \vec{X}_t) - \vec{X}_t \right) dt + \vec{\xi} \odot d\vec{W}_t^v\]

\emph{(Note: \(\odot\) denotes the element-wise Hadamard product).}

The correlated Wiener processes governing the asset returns and
volatility shocks are defined in terms of independent standard normal
variates \(\vec{Z}_1, \vec{Z}_2 \sim \mathcal{N}(\vec{0}, I_N)\)
transformed via Cholesky decomposition:

\[d\vec{Z}_1^{\text{corr}} = L \, d\vec{Z}_1\]

where \(L\) is the lower triangular Cholesky factor of the equity return
correlation matrix \(\Sigma\). The transformed shocks are used for both
asset returns and volatility:

\[d\vec{W}_t^S = L \, d\vec{Z}_1\]

\[d\vec{W}_t^v = \vec{\rho} \odot L \, d\vec{Z}_1 + \sqrt{1 - \vec{\rho}^2} \odot d\vec{Z}_2\]

\textbf{Parameter Definitions:}

\begin{itemize}
\tightlist
\item
  \(\vec{r}\): \(N \times 1\) risk-free rate vector.
\item
  \(\vec{q}\): \(N \times 1\) continuously compounded dividend yield
  vector.
\item
  \(\vec{\theta}\): \(N \times 1\) idiosyncratic baseline log-variance
  target vector.
\item
  \(\vec{\kappa}\): \(N \times 1\) mean-reversion speed vector.
\item
  \(\vec{\gamma}\): \(N \times 1\) network volatility sensitivity
  vector.
\item
  \(W\): \(N \times N\) directed network edge weight matrix (rows
  normalized to sum to 1).
\item
  \(\vec{\xi}\): \(N \times 1\) volatility-of-volatility vector.
\item
  \(\vec{\rho}\): \(N \times 1\) vector of asset-specific leverage
  correlations \(\in [-1, 1]\) (each element represents return-variance
  correlation for that asset).
\item
  \(\Sigma\): \(N \times N\) empirical correlation matrix of equity
  log-returns.
\item
  \(L\): \(N \times N\) lower triangular Cholesky decomposition matrix
  such that \(\Sigma = LL^T\).
\end{itemize}

\begin{center}\rule{0.5\linewidth}{0.5pt}\end{center}

\subsubsection{2. Calibration Pipeline}\label{calibration-pipeline}

\paragraph{2.1 Hybrid Approach: Structural Mask + Statistical
Weighting}\label{hybrid-approach-structural-mask-statistical-weighting}

The \passthrough{\lstinline!MarketDataFetcher/data\_fetcher.py!}
pipeline combines \textbf{domain knowledge} (user-provided network
structure) with \textbf{statistical validation} (Diebold-Yilmaz FEVD):

\textbf{Stage 1: Equity Return Correlations (Cholesky Decomposition)}

\begin{enumerate}
\def\labelenumi{\arabic{enumi}.}
\tightlist
\item
  \textbf{Fetch historical equity prices} from Yahoo Finance for the
  \(N\) assets over a lookback window (typically 252 trading days)
\item
  \textbf{Compute log-returns}: \(r_{i,t} = \ln(P_{i,t} / P_{i,t-1})\)
  for each asset \(i\)
\item
  \textbf{Calculate correlation matrix}:
  \(\Sigma = \text{Corr}(\{r_{i,t}\})\) (Pearson correlation of returns)
\item
  \textbf{Ensure positive semi-definiteness} via eigenvalue correction
  if numerical instabilities exist
\item
  \textbf{Compute Cholesky decomposition}:
  \(L = \text{cholesky}(\Sigma)\) such that \(\Sigma = LL^T\)
\item
  \textbf{Export to CSV}:
  \passthrough{\lstinline!cholesky\_L\_matrix.csv!} (lower triangular
  \(N \times N\) matrix)
\end{enumerate}

\textbf{Stage 2: Network Structure via Mask + Diebold-Yilmaz Weighting}

The framework uses a \textbf{two-step hybrid approach} to construct the
directed network matrix \(W\):

\textbf{Step 2a: User-Defined Structural Mask} (Binary, Non-stochastic)

User creates \passthrough{\lstinline!skeleton\_W\_matrix.csv!}
specifying which relationships are economically meaningful:

\begin{lstlisting}
From,To
A,B
A,C
B,C
\end{lstlisting}

\textbf{Step 2b: Statistical Weighting (Diebold-Yilmaz)}

\begin{enumerate}
\def\labelenumi{\arabic{enumi}.}
\tightlist
\item
  \textbf{Fetch historical implied volatility (IV) data} from DoltHub
  for the same \(N\) assets
\item
  \textbf{Fit a VAR(1) model} to the IV time series
\item
  \textbf{Calculate Generalized Forecast Error Variance Decomposition
  (GFEVD)} using the Diebold-Yilmaz methodology
\item
  \textbf{Merge with Structural Mask}: Combine user-defined binary
  relationships with Diebold-Yilmaz weights
\item
  \textbf{Enforce sparsity}: retain only top \(K\) neighbors per asset
  (typically \(K=3\)) to optimize computational efficiency
\item
  \textbf{Ensure row-stochasticity}: normalize each row to sum to 1.0
\item
  \textbf{Export to CSV}:
  \passthrough{\lstinline!calibrated\_W\_matrix.csv!} (row-normalized
  \(N \times N\) matrix)
\end{enumerate}

\paragraph{2.2 Data Separation
Rationale}\label{data-separation-rationale}

The framework uses \textbf{two distinct data sources} to model
complementary phenomena:

{\def\LTcaptype{none} % do not increment counter
\begin{longtable}[]{@{}
  >{\raggedright\arraybackslash}p{(\linewidth - 6\tabcolsep) * \real{0.1951}}
  >{\raggedright\arraybackslash}p{(\linewidth - 6\tabcolsep) * \real{0.3171}}
  >{\raggedright\arraybackslash}p{(\linewidth - 6\tabcolsep) * \real{0.2683}}
  >{\raggedright\arraybackslash}p{(\linewidth - 6\tabcolsep) * \real{0.2195}}@{}}
\toprule\noalign{}
\begin{minipage}[b]{\linewidth}\raggedright
Matrix
\end{minipage} & \begin{minipage}[b]{\linewidth}\raggedright
Data Source
\end{minipage} & \begin{minipage}[b]{\linewidth}\raggedright
Mechanism
\end{minipage} & \begin{minipage}[b]{\linewidth}\raggedright
Purpose
\end{minipage} \\
\midrule\noalign{}
\endhead
\bottomrule\noalign{}
\endlastfoot
\textbf{L} (Cholesky) & Equity prices (log-returns) & Direct
contemporaneous correlation & Captures fundamental business cycle
synchronization; affects immediate price co-movements \\
\textbf{W} (Network) & Implied volatility (IV) & Spillover amplification
& Captures volatility clustering and regimes; affects volatility
contagion across assets \\
\end{longtable}
}

\begin{center}\rule{0.5\linewidth}{0.5pt}\end{center}

\subsubsection{3. Discretization \& Path Simulation
Algorithm}\label{discretization-path-simulation-algorithm}

To generate discrete Monte Carlo paths, the continuous-time system is
approximated over \(T\) time steps using an Euler-Maruyama scheme for
the log-variance process and an exact exponential time-stepping solution
for the asset prices.

Given initial conditions \(\vec{S}_0\) and \(\vec{X}_0\), for each time
step \(t \in \{1, 2, \dots, T\}\):

\textbf{Step 1: Generate Independent Standard Normal Shocks}

Generate two independent standard normal vectors:
\[\vec{Z}_1, \vec{Z}_2 \sim \mathcal{N}(\vec{0}, I_N)\]

\textbf{Step 2: Apply Cholesky Transformation}

Transform the independent shocks via the Cholesky matrix to induce
multi-asset return correlations:
\[\vec{Z}_1^{\text{corr}} = L \, \vec{Z}_1\]

This vector is used for \textbf{both} asset return and volatility
processes.

\textbf{Step 3: Construct Brownian Increments}

Construct Brownian increments for asset returns:
\[\Delta \vec{W}_t^S = \sqrt{\Delta t} \, \vec{Z}_1^{\text{corr}}\]

Construct Brownian increments for volatility using the \textbf{same
transformed shocks} to preserve the Heston leverage effect:
\[\Delta \vec{W}_t^v = \vec{\rho} \odot \sqrt{\Delta t} \, \vec{Z}_1^{\text{corr}} + \sqrt{1 - \vec{\rho}^2} \odot \sqrt{\Delta t} \, \vec{Z}_2\]

\textbf{Key Design Principle}: Both processes use
\(\vec{Z}_1^{\text{corr}}\), ensuring that:

\begin{itemize}
\tightlist
\item
  \textbf{Multi-asset return correlations}: Encoded in the \(L\) matrix
\item
  \textbf{Heston leverage effect}: Preserved via correlation \(\rho\)
  within each asset
\item
  \textbf{Cross-asset volatility spillovers}: Modeled through the
  network matrix \(W\) in the variance dynamics
\end{itemize}

\textbf{Step 4: Compute Spatiotemporal Volatility Contagion}

Evaluate the network-weighted volatility spillover:
\[\vec{N}_t = W\vec{X}_{t-1}\]

Compute the dynamic mean-reversion target incorporating both
idiosyncratic and network-driven components:
\[\vec{\Theta}_t = \vec{\theta} + \vec{\gamma} \odot (\vec{N}_t - \vec{X}_{t-1})\]

\textbf{Step 5: Update Network Log-Variance (Euler-Maruyama)}

\[\vec{X}_t = \vec{X}_{t-1} + \vec{\kappa} \odot \left( \vec{\Theta}_t - \vec{X}_{t-1} \right) \Delta t + \vec{\xi} \odot \Delta \vec{W}_t^v\]

\textbf{Step 6: Transform to Real Variance}

\[\vec{V}_{t-1} = \exp(\vec{X}_{t-1})\]

Exponential transformation guarantees strict positivity:
\(\vec{V}_{t-1} > 0\).

\textbf{Step 7: Update Asset Prices (Exact Log-Normal Discretization)}

Construct the risk-neutral drift with dividend yield adjustment and Itô
correction:

\[\vec{\mu}_t = \left(\vec{r} - \vec{q} - \frac{1}{2}\vec{V}_{t-1}\right)\Delta t\]

Construct the diffusion component using correlated Brownian increments:

\[\vec{\sigma}_t = \sqrt{\vec{V}_{t-1}} \odot \Delta \vec{W}_t^S\]

Apply the log-normal update:

\[\vec{S}_t = \vec{S}_{t-1} \odot \exp(\vec{\mu}_t + \vec{\sigma}_t)\]

This exact discretization ensures all simulated prices remain strictly
positive and achieves zero structural drift for the asset GBM.

\begin{center}\rule{0.5\linewidth}{0.5pt}\end{center}

\subsubsection{4. Pricing Methods}\label{pricing-methods}

\paragraph{4.1 European Options}\label{european-options}

Simple discounted expectation over all Monte Carlo paths:

\[C_{\text{Euro}}^0 = e^{-r T} \mathbb{E}[C(S_T, K)]\]

\paragraph{4.2 American Options (Longstaff-Schwartz
LSMC)}\label{american-options-longstaff-schwartz-lsmc}

Backward induction using least squares regression on polynomial basis
functions that include network effects:

\[V_t(\vec{S}_t, \vec{X}_t) = \max\left\{ \text{Payoff}(S_t), \mathbb{E}[V_{t+1} \mid \vec{S}_t, \vec{X}_t] \right\}\]

The continuation value is projected onto a 4-term or 5-term basis:

\begin{itemize}
\tightlist
\item
  \textbf{Constant term}: 1
\item
  \textbf{Log-moneyness}: \(\ln(S/K)\)
\item
  \textbf{Convexity}: \([\ln(S/K)]^2 - 1\)
\item
  \textbf{Local volatility}: \(X_t\) (log-variance state)
\item
  \textbf{Network spillover} (if active): \(\sum_j W_{ij} X_{j,t}\)
  (weighted sum of neighbor variances)
\end{itemize}

\begin{center}\rule{0.5\linewidth}{0.5pt}\end{center}

\subsubsection{5. Workflow \& User Inputs}\label{workflow-user-inputs}

\textbf{Required User Inputs:}

\begin{enumerate}
\def\labelenumi{\arabic{enumi}.}
\tightlist
\item
  \textbf{Ticker list} (e.g.,
  \passthrough{\lstinline!["AAPL", "MSFT", "JPM", ...]!})
\item
  \textbf{Skeleton network mask}
  (\passthrough{\lstinline!skeleton\_W\_matrix.csv!})

  \begin{itemize}
  \tightlist
  \item
    Binary matrix specifying which relationships matter
  \item
    User provides domain knowledge (Bloomberg terminal, economic
    reasoning)
  \end{itemize}
\item
  \textbf{FEVD horizon} (typically 10-20 steps)

  \begin{itemize}
  \tightlist
  \item
    Longer horizon = stronger spillover effects
  \end{itemize}
\end{enumerate}

\begin{center}\rule{0.5\linewidth}{0.5pt}\end{center}

\subsubsection{6. Advantages of Hybrid
Approach}\label{advantages-of-hybrid-approach}

\begin{itemize}
\tightlist
\item
  \textbf{Domain-aware}: Users encode domain knowledge directly (network
  structure)
\item
  \textbf{Statistically validated}: Diebold-Yilmaz provides empirical
  spillover strengths
\item
  \textbf{Parsimonious}: Avoids spurious spillovers by masking
  irrelevant pairs
\item
  \textbf{Flexible}: Skeleton can be updated based on new Bloomberg
  terminal data
\item
  \textbf{Interpretable}: Each edge represents an economically
  meaningful relationship
\item
  \textbf{Robust}: Doesn't require dense historical data for all
  possible pairs
\end{itemize}

\begin{center}\rule{0.5\linewidth}{0.5pt}\end{center}

\subsubsection{7. Key Implementation
Features}\label{key-implementation-features}

\begin{itemize}
\tightlist
\item
  \textbf{Log-Normal Exact Discretization}: Eliminates bias drift for
  underlying asset paths
\item
  \textbf{Eigenvalue Correction}: Ensures correlation matrix positive
  semi-definiteness before Cholesky decomposition
\item
  \textbf{Col-Pivot QR for Regression}: Robust least squares solution
  even when basis functions are nearly collinear
\item
  \textbf{Adaptive Basis Functions}: LSMC automatically includes network
  spillover term only when asset has active neighbors
\item
  \textbf{Row-Stochastic W Matrix}: Ensures economic interpretability of
  volatility spillovers
\item
  \textbf{Unified Shock Space}: Both asset returns and volatility
  leverage use Cholesky-transformed shocks, preserving the Heston
  framework
\end{itemize}

\end{document}
